% !TEX root = ./main.tex
\chapter{Introduzione}\label{chapintroduzione}

In molti ambienti di realtà virtuale (Virtual Reality, VR), che prevedono l'interazione con un avatar, la scena regge visivamente, ma la conversazione con l'agente resta fragile o troppo rigida; allo stesso tempo l'identità del personaggio tende a disperdersi tra una sessione e l'altra. SOULFRAME nasce da questo scarto: il prototipo permette di creare o ricostruire un avatar, dargli un aspetto riconoscibile, associargli una voce e parlargli in una scena 3D senza ripartire da zero a ogni incontro. Il codice sorgente del progetto è disponibile nella repository dedicata \cite{SOULFRAMERepo}.

La difficoltà nel creare un agente conversazionale realistico sta nel tenere insieme tre aspetti che spesso vengono affrontati separatamente: interazione vocale a turni, continuità del dialogo e presenza dell'interlocutore virtuale. In VR non basta che l'avatar ``ci sia'': la percezione di presenza dipende anche dalla coerenza tra azioni, tempi di risposta ed eventi della scena \cite{Slater2016VR}. Quando questi elementi si disallineano, l'agente appare poco affidabile anche se il contenuto linguistico è corretto.

Le evidenze sperimentali sugli agenti conversazionali embodied (Embodied Conversational Agents, ECA)\footnote{Nel testo uso l'acronimo ECA nel senso impiegato dalla letteratura sugli agenti conversazionali dotati di una rappresentazione corporea; il quadro teorico viene approfondito nel Capitolo~\ref{chap:stato-arte}.} indicano che una pipeline vocale in tempo reale basata su Speech-to-Text (STT) e Text-to-Speech (TTS) può aumentare la co-presenza rispetto a soluzioni preregistrate \cite{Kan2023ECA_Training}. La letteratura in realtà estesa (Extended Reality, XR) mostra anche che molti prototipi sono sviluppati in VR, spesso con Unity e con visore Head-Mounted Display (HMD), e che cresce l'interesse per dialoghi più adattivi supportati da modelli neurali \cite{Yang2025ECA_XR}. SOULFRAME si muove in questo spazio, ma in una configurazione senza HMD obbligatorio e centrata su un client Unity.

Alla base del progetto c'è l'idea che l'identità dell'agente dipenda dalla tenuta congiunta di volto, voce e tracce lasciate nelle interazioni precedenti. Da qui nasce la domanda se questo insieme renda l'interlocutore virtuale più credibile di un semplice chatbot collocato in scena. Nel Capitolo~\ref{chap:stato-arte}, Figura~\ref{fig:ca-evoluzione}, richiamo la traiettoria dei Conversational Agents dalla fase scriptata alle soluzioni recenti basate su modelli linguistici di grandi dimensioni (Large Language Models, LLM).

\section{Perimetro e requisiti di progetto}\label{secperimetro-requisiti}
SOULFRAME è un prototipo di interazione immersiva con ECA vocale embodied in ambiente 3D. L'obiettivo è valutarne fattibilità tecnica e qualità percepita in scenari con componente sociale, mantenendo riconoscibile nel tempo il personaggio costruito a partire dai dati dell'utente. Restano fuori una piattaforma generale per creare mondi virtuali e una soluzione clinica certificata.

L'interazione resta orientata a una sola persona per sessione: non implementa sincronizzazione multiutente né gestione concorrente di più interlocutori sullo stesso contesto conversazionale.

Sul fondo c'è anche il tema della solitudine e delle relazioni mediate. In questa tesi, però, il contributo resta tecnico: costruire una base conversazionale embodied abbastanza stabile su quel dominio, senza attribuire al prototipo effetti terapeutici o risultati che non potrebbe sostenere.

Il lavoro si concentra sull'integrazione \textit{end-to-end} tra componente vocale, memoria conversazionale e presenza dell'agente nella scena. Per questo la conversazione vera e propria si apre solo dopo un onboarding minimo del profilo: senza un riferimento vocale e senza contenuti recuperabili, l'interazione tende a diventare quasi \textit{stateless}.

La conversazione segue una pipeline a turni di tipo \textit{push-to-talk} (premi-per-parlare): acquisizione audio, trascrizione STT, orchestrazione conversazionale con memoria e risposta TTS. Nel flusso operativo il turno può passare anche da input testuale, ma resta legato alla stessa sessione. L'LLM viene usato per generare la risposta a partire dal contesto del turno, non come archivio personale del personaggio: i contenuti rilevanti vengono recuperati da una memoria esterna e poi passati al modello. Descrizione immagini e OCR, cioè riconoscimento ottico del testo, ampliano il materiale richiamabile durante la conversazione.

I requisiti non funzionali riguardano soprattutto latenza percepita, robustezza ai guasti parziali e modularità dei servizi. La tesi adotta quindi criteri operativi verificabili sul prototipo, evitando soglie astratte non allineate al contesto d'uso reale.

\section{Obiettivi e contributi di SOULFRAME}\label{secobiettivi-contributi}
Nel filone degli ECA in XR, dove sono comuni implementazioni VR in Unity e interazioni vocali ma i metodi di valutazione restano eterogenei \cite{Yang2025ECA_XR}, SOULFRAME prova a rendere osservabile l'integrazione \textit{end-to-end} con criteri espliciti, senza separare la qualità del dialogo dalla persistenza del personaggio.

\textbf{RQ1} riguarda la fattibilità tecnica di un'interazione vocale in tempo reale con un agente embodied in ambiente immersivo 3D. In termini operativi, la catena STT $\rightarrow$ orchestrazione conversazionale con memoria $\rightarrow$ TTS $\rightarrow$ output audiovisivo deve funzionare in modo continuo, includendo anche le capacità di supporto (descrizione immagini e OCR documentale). La valutazione considera tempi per blocco, completamento dei turni senza errori bloccanti e coerenza delle transizioni tra ascolto, elaborazione e risposta.

\textbf{RQ2} riguarda la qualità percepita dell'interazione, con attenzione a co-presenza, naturalezza dialogica e differenze introdotte dalla personalizzazione dell'avatar. In questo lavoro la parte empirica confronta infatti una condizione con avatar neutro e una con avatar personale, osservando soprattutto fiducia nell'agente, percezione dell'avatar e variazione dello stato affettivo. Qualità percepita e co-presenza non vengono usate come proxy di riduzione della loneliness: servono piuttosto a capire se l'infrastruttura embodied regge abbastanza da rendere sensati studi successivi su quel piano. Studi comparativi tra agenti conversazionali e audio pre-scriptato mostrano differenze osservabili, soprattutto sulla co-presenza \cite{Kan2023ECA_Training}.

\textbf{RQ3} riguarda modularità e riproducibilità della pipeline. La pipeline deve permettere la sostituzione dei componenti principali senza riscrivere l'intera architettura e rendere tracciabili configurazioni, tempi e risultati delle prove; è un requisito strutturale dell'obiettivo generale, perché rende più solida la portabilità del prototipo tra ambienti diversi. Figura~\ref{fig:obiettivi-contributi} sintetizza questa corrispondenza.

\begin{figure}[t]
\centering
\includegraphics[width=1.0\textwidth]{obiettivi_contributi.png}
\caption[Obiettivi e contributi di SOULFRAME]{Corrispondenza tra obiettivi di ricerca e contributi del progetto SOULFRAME.}
\label{fig:obiettivi-contributi}
\end{figure}

Sul piano tecnico, SOULFRAME adotta un'architettura modulare \textit{end-to-end} basata su micro-servizi con interfacce applicative esplicite (Application Programming Interface, API). La verifica assume la forma di un \textit{proof-of-concept} osservabile con un set minimo di misure soggettive e osservazioni qualitative. La documentazione delle scelte progettuali sostiene anche replicabilità e trasferibilità del sistema, in linea con il dibattito sul \textit{dialogue management} e sulla necessità di metriche più comparabili \cite{Laranjo2018CA}.

\section{Panoramica del sistema e flusso end-to-end}\label{secpanoramica-sistema}
SOULFRAME adotta un'architettura client--server per sostenere un dialogo vocale a turni in ambiente immersivo 3D. Il client principale è implementato in Unity: la modalità principale di distribuzione è la build WebGL (Web Graphics Library), affiancata durante sviluppo, verifica e confronto da una build Windows nativa dello stesso frontend; il backend gestisce trascrizione, recupero del contesto, generazione testuale e sintesi vocale. Sul lato utente, il flusso va dalla selezione o creazione dell'avatar alla preparazione del profilo, fino alla conversazione vera e propria.

Il flusso può apparire lineare: input dell'utente, trascrizione STT, costruzione della risposta contestuale, sintesi TTS e playback sull'avatar. Nel sistema reale, però, il blocco conversazionale interpreta il turno, instrada la richiesta e combina contesto recente, contenuti recuperati e risposta finale. Tutte queste fasi devono restare coordinate nel tempo, perché ritardi o transizioni incoerenti compromettono la naturalezza percepita.

Si è scelto di distribuire il client principalmente in WebGL, senza HMD obbligatorio, per abbassare la barriera di accesso durante sviluppo e test. La build Windows non rappresenta quindi un prodotto separato, ma un termine di confronto utile per isolare gli effetti del runtime: la logica applicativa resta la stessa, mentre cambiano vincoli del browser, continuità del playback audio e modalità di distribuzione.

La vista d'insieme architetturale è discussa in \autoref{sec:architettura-soulframe} e sintetizzata in Figura~\ref{fig:pipeline-e2e}, dove la pipeline viene ripresa con maggior dettaglio tecnico.\label{secmetodo-struttura}

\begin{figure}[t]
\centering
\includegraphics[width=0.95\textwidth]{struttura_tesi.png}
\caption[Struttura della tesi]{Struttura della tesi e organizzazione dei capitoli.}
\label{fig:struttura-tesi}
\end{figure}

Lo sviluppo di SOULFRAME è proceduto per prototipi incrementali. Era l'approccio più adatto a un sistema che combina reattività a turni (\textit{push-to-talk}), elaborazione linguistica e rendering immersivo: dapprima sono stati validati i moduli in isolamento, poi sono stati integrati nella pipeline completa e infine verificati in scenari via via più complessi, mantenendo tracciabili compromessi e dipendenze.

La tesi segue la stessa logica: dai fondamenti teorici si passa all'architettura, poi all'implementazione, alla valutazione tecnica del prototipo e infine allo studio empirico sugli utenti.

Il Capitolo 2 ricostruisce stato dell'arte e quadro teorico; il Capitolo 3 descrive architettura, componenti e scelte implementative; il Capitolo 4 raccoglie i risultati tecnici su pipeline, latenza e casi di guasto; il Capitolo 5 introduce lo studio empirico comparativo tra condizione neutra e condizione personale; il Capitolo 6 chiude il lavoro tirando insieme limiti, contributi e sviluppi futuri.
